\section{Evidence at deep nesting: block pairs}\label{sec:dnp}

The deep-nesting pairs of the eighth paper of the series~\cite[Section~2.5]{BoundedNesting2026} are the sharpest open
instance of $\TCP(2)$: its Conjecture~A$''$~\cite[Conjecture~5.24]{BoundedNesting2026}, that at a suitable odd shift every
label of cost at most $n_{\mathrm h}/8$ on every cover doubly $\theta_0$-serves at most $C\log_2q+C'+\eta(m+1)$ levels with
$\eta<1$, would make them depth-robust, so that $\TCP(2)$ fails~\cite[Corollary~5.26]{BoundedNesting2026}. For pairs of
block labels of bounded length, the tenth paper~\cite[Section~5]{DirectedReciprocity2026} proved failure witnesses in three
regimes (its Theorem~A, which rests on its Theorems~V$^*$ and~LU, all Lean-verified there) and named four further steps
towards Conjecture~A$''$~\cite[Remark~5.6]{DirectedReciprocity2026}: (a)~one odd shift serving both directions;
(b)~half-blindness of the fine levels at that shift; (c)~from witnesses to counts of served levels; (d)~unbounded block
lengths and general labels.

This section closes steps~(b) and~(c) for bounded block length and reduces step~(a) to Conjecture~S (\cref{conj:S}): the
straddle counts of most coarse classes concentrate as the shift varies. The result is Theorem~BP (\cref{thm:BP}): assuming
Conjecture~S, Conjecture~A$''$ holds for pairs of block labels of length at most~$a_0$, at one odd shift, with no
$\log q$ term and no cost hypothesis. Theorem~BP is evidence that the deep-nesting pairs are depth-robust; it does not
decide $\TCP(2)$, since Conjecture~A$''$ concerns every label of cost at most $n_{\mathrm h}/8$ (\cref{rem:dnp-scope}).
\Cref{sec:dnp-partial} gives partial results on Conjecture~S and numerical evidence for it up to $m=4095$.

The written proofs use the tenth paper's Theorems~V$^*$, LU and~A (Lean-verified there) and the eighth paper's lemmas on
the deep-nesting pair (written proofs there). Every numbered result of this section is formalized in Lean, with
Conjecture~S as a hypothesis where the statement assumes it. The formalization includes the deep-nesting walk, the
eighth paper's Lemma~5.1(a)--(c) and its blindness bound (Proposition~5.17, without its level-$0$ term), so no step of the
formal proofs rests on a written proof of the eighth paper (\cref{sec:dnp-setting}).

\subsection{Setting}\label{sec:dnp-setting}

We use the notation of~\cite[Sections~2.5 and~5]{BoundedNesting2026}: $s_0>m\ge1$, $\Lambda=2^{s_0}$, $R_i=2^{s_0-i}$,
$n_i=m-i+1$, $\Pi_i=n_iR_i$ (the length of a stroke of level~$i$), $R_{m+1}=0$, and we put $N=m+1$, so $P=4N\Lambda=4\Pi_0$.
The pair is $x_1=U=U(s_0,m,1)$ and $x_2=x_1(\cdot+\xi)$ for an odd shift~$\xi$, so that no time is a turn of both heads.
On the $r$-fold cover a label~$G$ with $q$ values gives the hop $(p,u)$ the error $e(p,u)=(G(u)-G(p)-1)\bmod q$; the hops of
head~$h$ of level~$k$ form the class $(h,k)$; and $G$ \emph{$\theta$-serves} level~$k$ for head~$h$ if the class has total
error at most $\theta qn_{\mathrm h}/(2N)$, and \emph{doubly $\theta$-serves} it if it does so for both heads. Every class
has exactly $n_{\mathrm h}/(2N)=4r\Lambda$ hops~\cite[Lemma~5.1(c)]{BoundedNesting2026}, so $\theta$-service means mean
error at most $\theta q$. A hop is \emph{blind} if its half-open interval $(p,u]$ contains no turn of the other head.

\paragraph{Block pairs.} For $a\ge1$ and $K\in\Z$, $B(a,K)$ is the block label of~\cite[Definition~5.11]{BoundedNesting2026};
that definition takes $K\ge0$, but it and the identities of~\cite[Lemma~5.12]{BoundedNesting2026} used below hold verbatim
for every integer~$K$, as in the tenth paper. Its \emph{cancellers} are the levels~$l$ with $m-l\equiv a\pmod{a+1}$; it
has level parameters $A_l,E_l$, block mass $Q(a,K)=(2^a-1)(2K+2)$ and path-sum increments $c_\sigb(l)=1+A_l$ and
$c_\sigf(l)=2+A_l+E_l$, with $c_\sigf(\kappa)-c_\sigb(\kappa)=-Q(a,K)-1$ at every canceller~$\kappa$. For $q\mid r$ the
\emph{block-pair label} of $(B(a_1,K_1),B(a_2,K_2))$ modulo~$q$ on the $r$-fold cover is $G=(\Phi_1+\Phi_2)\bmod q$, where
$\Phi_h$ is the potential of head~$h$ carrying $B(a_h,K_h)$ at its own turns~\cite[Section~5.3]{BoundedNesting2026}; it is
well defined because each $\Phi_h$ grows by an integer per period and $q\mid r$, as
in~\cite[Proposition~5.23]{BoundedNesting2026}, and it is turn-supported. For a head~$h$, a level~$k$ and $j\ge0$ with
$k+j\le m$, the \emph{blind errors} of $B(a_h,K_h)$ are
\[
  a^{(h)}(k)=A^{(h)}_k\bmod q,\qquad
  b^{(h)}_j(k)=\Bigl(E^{(h)}_k-\sum_{i=k+1}^{k+j}\bigl(A^{(h)}_i+E^{(h)}_i+2\bigr)\Bigr)\bmod q
\]
and $\tilde Q_h=Q(a_h,K_h\bmod q)$; in the tenth paper's library these errors are
\lean{DirectedReciprocity.blindApexErr} and \lean{DirectedReciprocity.blindBandErr}. By the band
decomposition~\cite[Lemma~5.10]{BoundedNesting2026}, the apex hop $(\alpha-d,\alpha+d)$, $1\le d<R_k$, of a level-$k$~V of
head~$h$ has error $(A^{(h)}_k+F)\bmod q$, and its valley hop of band~$j$, at a cell $v+\sigma x$ with
$R_{k+j+1}<x<R_{k+j}$, has error $(b^{(h)}_j(k)+F)\bmod q$, where $F$ is the other head's increment over the hop; on a
blind hop $F=0$.

\paragraph{Four results of the tenth paper.} All four are Lean-verified there; the second and third are the shift-free
components \lean{DirectedReciprocity.block_overflow_band_witness} and \lean{DirectedReciprocity.block_straddle_residue_ge_half_q} of the proof
of~\cite[Theorem~5.5]{DirectedReciprocity2026} (its Regimes~2 and~3). Define $a(k)$, $b_j(k)$ as above for any integer
sequences $(A_l)$, $(E_l)$ and $q\ge3$.
\begin{enumerate}[label=(\Roman*)]
\item \emph{Run bound}~\cite[Theorem~5.3(2)]{DirectedReciprocity2026}. If $s\ge1$ and the levels $k_0,\dots,k_0+s-1$
  satisfy $8a(k_0+i)\le q$ and $8b_0(k_0+i)\le q$ for $i<s$, $8b_1(k_0+i)\le q$ for $i+1<s$, and $8b_2(k_0+i)\le q$ for
  $i+2<s$, then $s\le2\lfloor\log_2q\rfloor+2$.
\item \emph{Overflow bound.} If $a\ge1$, $K\in\Z$, $q\ge32$ and $4Q(a,K\bmod q)+4>q$, then there is
  $\mathrm{rem}\in\{0,\dots,a-1\}$ such that every level $k<m$ with $(m-k)\bmod(a+1)=\mathrm{rem}$ has
  $16a(k)\ge q$, $8b_0(k)\ge q$ or $8b_1(k)\ge q$, for the blind errors of $B(a,K)$.
\item \emph{Residue bound.} If $a_1,a_2\ge1$, $K_1,K_2\in\Z$, $q\ge1$, $\tilde Q_1\le\tilde Q_2$, $4\tilde Q_2+4\le q$, and
  $\kappa$ is a canceller of $B(a_2,\cdot)$, then for every level~$k$,
  $\bigl(2A^{(1)}_k+2(c^{(2)}_\sigf(\kappa)-c^{(2)}_\sigb(\kappa))\bigr)\bmod q\ge q/2$, where $A^{(1)}$ belongs to
  $B(a_1,K_1)$ and $c^{(2)}_\sigb,c^{(2)}_\sigf$ to $B(a_2,K_2)$.
\item \emph{Theorem~V$^*$}~\cite[Theorem~5.2(2)]{DirectedReciprocity2026}. Let $\Phi\colon\Z\to\Z$, let $k\ge1$
  and~$\kappa$ satisfy $\Pi_\kappa\le R_k$, let $\alpha\in\Z$, and suppose that there are integers $C_V,c_\sigb,c_\sigf$
  and a function~$Y$ with $\Phi(\alpha+o)=C_V+c_\sigb+Y(o)$ and $\Phi(\alpha+\Pi_\kappa+o)=C_V+c_\sigf+Y(o)$ for
  $0\le o<\Pi_\kappa$. For an integer~$c_0$ let $S=\mathrm{Strad}(c_0,\alpha,\Pi_\kappa,R_k)$ be the set of $d\in[1,R_k-1]$
  with $\alpha\le c_0-d<\alpha+\Pi_\kappa\le c_0+d<\alpha+2\Pi_\kappa$, and for $A'\in\Z$ and $q\ge1$ put
  $e(d)=(A'+\Phi(c_0+d)-\Phi(c_0-d))\bmod q$. Then
  $\card S\cdot\bigl((2A'+2(c_\sigf-c_\sigb))\bmod q\bigr)\le2\sum_{d\in S}e(d)$.
\end{enumerate}

\paragraph{Formalization.} The walk~$U$ is \lean{MergeDepth.dnpWalk}. Its turns are the apex and valley times of its V's
and two tips per period (\lean{MergeDepth.isTurnZ_dnpWalk_iff}), and the elements of the tenth paper's apex sets, used
for~$\mathcal A_i$ in \cref{def:dnp-straddle}, are apices of V's of level~$i$ (\lean{MergeDepth.dnpWalk_V}). The eighth
paper's Lemma~5.1 is proved: each head has $2^{m+3}-6$ turns per period (\lean{MergeDepth.card_turns_period}), the hops of
level $k\ge1$ are those listed there (\lean{MergeDepth.levelHops_eq_listed}), and every class has $4\Lambda$ hops per
period and $4r\Lambda=n_{\mathrm h}/(2N)$ on the $r$-fold cover (\lean{MergeDepth.levelCount_eq},
\lean{MergeDepth.card_coverClass}). A block label is a jump function with the properties
of~\cite[Definition~5.11]{BoundedNesting2026} (\lean{MergeDepth.IsBlockLabel}); there is exactly one, and its potential is
that of the tenth paper's stroke recursion (\lean{MergeDepth.IsBlockLabel.potentialOf_eq}). The block-pair label, its
errors, $\theta$-service and blindness are \lean{MergeDepth.blockPairLabel}, \lean{MergeDepth.hopError},
\lean{MergeDepth.ThetaServes} and \lean{MergeDepth.IsBlind}. \Cref{lem:dnp-family,lem:dnp-halfblind,lem:dnp-witness} and
Theorem~BP are first proved in one period, for the families of \cref{def:dnp-family} listed explicitly as pairs of times,
with these potentials, and with service measured on the four families of a class. They transfer to the walk pair because
the families consist of hops of their classes (\lean{MergeDepth.family_inHopClass}), the errors are those of the
block-pair label (\lean{MergeDepth.hopError_blockPairLabel}), and a level that is $\theta$-served for a head on some cover
is also served when service is measured on its four families (\lean{MergeDepth.ThetaServes.servedOnWitness}). The formal
proofs do not use $q\mid r$, which is needed only for $G$ to be a label of the cover
(\lean{MergeDepth.blockPairLabel_add_mul}).

\subsection{Straddle counts and Conjecture S}\label{sec:dnp-S}

\begin{definition}[Straddle counts and good levels]\label{def:dnp-straddle}
\formal{MergeDepth/BlockPairsProfile.lean}{MergeDepth.coarseZone,MergeDepth.jZero,MergeDepth.goodLevels,MergeDepth.goodCount}
Let $\mathcal K=\{1,\dots,\lfloor N/4\rfloor\}$ be the \emph{coarse zone} and $j_0=\lfloor\log_2N\rfloor+1$. For $a\ge1$
and $k\ge1$ let $\kappa_a(k)$ be the least integer $\kappa\ge k+j_0$ with $m-\kappa\equiv a\pmod{a+1}$ (the tenth paper's
\lean{DirectedReciprocity.cancellerBelow}$(m,a,k)$ if $k+j_0\le m$); thus $j_0\le\kappa_a(k)-k\le j_0+a$, and $\kappa_a(k)$ is a canceller of
$B(a,\cdot)$ if $\kappa_a(k)\le m$. For $1\le i\le m$ let $\mathcal A_i\subseteq[0,P)$ be
the set of the $2^{i+1}$ apex times of the level-$i$ V's of~$U$ in one period. For $1\le k<\kappa\le m$ with
$\Pi_\kappa\le R_k$ and $\xi\in\Z$ put
\[
  X^\pm_{k,\kappa}(\xi)=\sum_{\alpha_1\in\mathcal A_k}\ \sum_{\alpha\in\mathcal A_\kappa}
  \card{\mathrm{Strad}\bigl((\alpha_1\pm\xi)\bmod P,\ \alpha,\ \Pi_\kappa,\ R_k\bigr)},
\]
which for $0\le\xi\le P$ is the tenth paper's \lean{DirectedReciprocity.classStraddleCount} at
$(s_0,m,k,\kappa,\mathrm{dir},\xi)$, with $\mathrm{dir}$ true for~$+$ and false for~$-$. A level $k\in\mathcal K$ is
\emph{$a$-good in direction~$\pm$} at~$\xi$ if
$X^\pm_{k,\kappa_a(k)}(\xi)\ge2^{k+2}R_k/2^{a+7}$, and $\Gamma_{a,\pm}(\xi)$ is the number of such levels.
\end{definition}

In head~$2$'s own time, where it is~$U$, the apex hop $(\alpha_1-d,\alpha_1+d)$ of head~$1$ is the window of
half-width~$d$ around $\alpha_1+\xi$; so $X^+_{k,\kappa}(\xi)$ counts the apex hops with $d<R_k$ of the class $(1,k)$ that
start in the back stroke and end in the forward stroke of a level-$\kappa$~V of head~$2$, and $X^-$ does the same with the
heads exchanged. The threshold is that of~\cite[Theorem~5.5(3)]{DirectedReciprocity2026}. Given $a_0\ge1$, we use the
\emph{size condition} $4(j_0+a_0)\le N$\quad\textup{($\ast$)}; under it, for $k\in\mathcal K$, $1\le a\le a_0$ and
$\kappa=\kappa_a(k)$, the tenth paper's \lean{DirectedReciprocity.cancellerBelow_props} gives $\kappa\le m$, $\Pi_\kappa\le R_k$ and
$N\le2n_\kappa$.

\begin{definition}[Straddle profile]\label{def:dnp-profile}
\formal{MergeDepth/BlockPairsProfile.lean}{MergeDepth.piHat,MergeDepth.apexHat,MergeDepth.valleyHat,MergeDepth.tent,MergeDepth.perTent,MergeDepth.profile,MergeDepth.profileMean,MergeDepth.profileVariance,MergeDepth.relVariance,MergeDepth.relVarianceA,MergeDepth.meanRelVariance}
Let $1\le k<\kappa\le m$ and $\hat\Pi_i=n_i2^{-i}$, so that $\Pi_i=\Lambda\hat\Pi_i$. For $S\subseteq\{1,\dots,k-1\}$ and
$S'\subseteq\{1,\dots,\kappa-1\}$ put
\[
  \hat a_S=\sum_{i=0}^{k-1}2^{-i}+\sum_{i\in S}\hat\Pi_i,\qquad
  \hat v_{S'}=\sum_{i=0}^{\kappa-1}2^{-i}+\hat\Pi_\kappa+\sum_{i\in S'}\hat\Pi_i,
\]
let $\Delta(y)=\max(0,\hat\Pi_\kappa-2\abs y)$, and define the $N$-periodic function
\[
  x_k(t)=\frac1{2^{k-1}\hat\Pi_\kappa}\sum_{S}\ \sum_{S'}\ \sum_{l\in\Z}\Delta\bigl(\hat a_S+t-\hat v_{S'}-lN\bigr)
  \qquad(t\in\R).
\]
For $t$ uniform on $[0,N)$ let $\mu_k=\mathbb E\,x_k(t)$ and $V_k=\operatorname{Var}x_k(t)/\mu_k^2$. Write $V^{(a)}_k$ for
$V_k$ at $\kappa=\kappa_a(k)$, which is at most~$m$ for $k\in\mathcal K$ once $4(j_0+a)\le N$, and
$\bar V^{(a)}_m=\card{\mathcal K}^{-1}\sum_{k\in\mathcal K}V^{(a)}_k$.
\end{definition}

These depend on $m$, $k$ and~$\kappa$, not on~$s_0$. In units of~$\Lambda$, the $\hat a_S$ and $\hat v_{S'}$ are the
level-$k$ apex and level-$\kappa$ valley times of one level-$0$ stroke, and $x_k$ is the normalized straddle count at a real
shift (\cref{lem:dnp-transfer}). Each tent has integral $\hat\Pi_\kappa^2/2$, so $\mu_k=n_\kappa/(4N)$.

\begin{conjecture}[Conjecture S]\label{conj:S}
For every integer $a\ge1$, $\bar V^{(a)}_m\to0$ as $m\to\infty$.
\end{conjecture}

In Lean, Conjecture~S for one~$a$ is \lean{MergeDepth.ConjectureS}.

By \cref{prop:S-fourier}, Conjecture~S is an averaged $L^2$ bound for Riesz products along the mixed-radix digits
$n_i2^{-i}$; by \cref{prop:S-resonance}, its analogue for single classes fails.

\begin{lemma}[Straddle counts and the profile]\label{lem:dnp-transfer}
\formal{MergeDepth/BlockPairsTransfer.lean}{MergeDepth.prob_straddle_vs_profile}
Let $1\le k<\kappa\le m<s_0$ with $\Pi_\kappa\le R_k$, and put $M_k=2^{k+1}\Pi_\kappa$.
\begin{enumerate}[label=(\alph*)]
\item For all integers $c_0$ and $\alpha$,
  $\card{\mathrm{Strad}(c_0,\alpha,\Pi_\kappa,R_k)}=\max\bigl(0,\Pi_\kappa-\abs{2(c_0-\alpha-\Pi_\kappa)+1}\bigr)$.
\item For every integer $0\le\xi\le P$, $X^+_{k,\kappa}(\xi)=M_k\,x_k\bigl((\xi+\frac12)/\Lambda\bigr)$ and
  $X^-_{k,\kappa}(\xi)=X^+_{k,\kappa}(P-\xi)$.
\item If $\xi$ is uniform over the odd residues modulo~$P$, then for every real~$\tau$ and each sign
  \[
    \Bigl|\Pr\bigl[X^\pm_{k,\kappa}(\xi)<\tau M_k\bigr]-\frac1N\,\mathrm{Leb}\{t\in[0,N):x_k(t)<\tau\}\Bigr|
    \ \le\ 2^{\kappa+2-s_0}.
  \]
\end{enumerate}
\end{lemma}

\begin{outline}
(a) With $z=c_0-\alpha$ the conditions say $\max(1,z-\Pi_\kappa+1,\Pi_\kappa-z)\le d\le\min(R_k-1,z,2\Pi_\kappa-1-z)$.
The integers $z-\Pi_\kappa+1$, $\Pi_\kappa-z$ sum to~$1$ and $z$, $2\Pi_\kappa-1-z$ sum to $2\Pi_\kappa-1$, so the bounds
$1$ and $R_k-1\ge\Pi_\kappa-1$ never bind; with $y=z-\Pi_\kappa$ the remaining bounds are $\frac12+\abs{y+\frac12}$ and
$\Pi_\kappa-\frac12-\abs{y+\frac12}$, leaving $\max(0,\Pi_\kappa-\abs{2y+1})$ values of~$d$.

(b) By~(a), $\card{\mathrm{Strad}(c_0,\alpha,\Pi_\kappa,R_k)}=\Lambda\Delta((c_0+\frac12-\nu)/\Lambda)$ with
$\nu=\alpha+\Pi_\kappa$ the valley time, and it vanishes unless $c_0\in[\alpha,\alpha+2\Pi_\kappa)$. The four level-$0$
strokes of a period have length $\Pi_0=N\Lambda$ and the same turn times relative to their starts, and every~V lies inside
one of them~\cite[Section~2.5]{BoundedNesting2026}; a stroke of level $l<m$ starting at~$t_0$ has the apex of its~V at
$t_0+R_l$, and the strokes of that~V start at $t_0+R_l$ and $t_0+R_l+\Pi_{l+1}$. So in the stroke starting at~$0$ the
level-$k$ apex times are the $\Lambda\hat a_S$ and the level-$\kappa$ valley times are the $\Lambda\hat v_{S'}$, and for
integers $c_0\in[0,P)$,
$\sum_{\alpha\in\mathcal A_\kappa}\card{\mathrm{Strad}(c_0,\alpha,\Pi_\kappa,R_k)}=\Lambda H((c_0+\frac12)/\Lambda)$ with the
$N$-periodic $H(y)=\sum_{S'}\sum_{l\in\Z}\Delta(y-\hat v_{S'}-lN)$. As $\mathcal A_k=\{\Lambda\hat a_S+j\Pi_0:0\le j\le3\}$,
$X^+_{k,\kappa}(\xi)=4\Lambda\sum_SH(\hat a_S+t)=M_kx_k(t)$ with $t=(\xi+\frac12)/\Lambda$. The second identity holds as
$(\alpha_1-\xi)\bmod P=(\alpha_1+(P-\xi))\bmod P$.

(c) Let $g(u)=M_kx_k(u/\Lambda)$, so $X^+_{k,\kappa}(\xi)=g(\xi+\frac12)$ by~(b). As the $\hat a_S,\hat v_{S'}$ are
multiples of $2^{-\kappa}$ and $\hat\Pi_\kappa/2$ is a multiple of $2^{-\kappa-1}$, $g$ is affine on each of the $N2^{\kappa+3}$
intervals $[iR_\kappa/2,(i+1)R_\kappa/2)$ that tile $[0,P)$. So each of them meets $\{g<\tau M_k\}$ in an interval of some
length~$\lambda$, which contains $\lambda/2\pm1$ of the points $\xi+\frac12$ with $\xi$ odd, as these have spacing~$2$.
Summing over the intervals, the number of odd $\xi\in[0,P)$ with $X^+_{k,\kappa}(\xi)<\tau M_k$ differs from half the measure
of $\{u\in[0,P):g(u)<\tau M_k\}$, which is $4\Lambda\,\mathrm{Leb}\{t\in[0,N):x_k(t)<\tau\}$ by periodicity, by at most
$N2^{\kappa+3}$; dividing by the number $P/2=2N\Lambda$ of odd residues gives the bound. The sign~$-$ follows from~(b), as
$\xi\mapsto P-\xi$ permutes the odd residues.
Formally, (a) is \lean{MergeDepth.card_straddleSet_eq}, which needs only $\Pi_\kappa\le R_k$; (b) is
\lean{MergeDepth.classStraddleCount_eq_profile} and \lean{MergeDepth.classStraddleCount_false_eq}; and~(c) is
\lean{MergeDepth.prob_straddle_vs_profile}.
\end{outline}

\begin{proposition}[Conjecture S yields a good shift]\label{prop:S-shift}
\formal{MergeDepth/BlockPairsShift.lean}{MergeDepth.notGoodProb_le,MergeDepth.prob_some_bad_le,MergeDepth.good_shift_prob}
Let $a_0\ge1$ and $0\le\gamma<1$, let $m<s_0$ satisfy~\textup{($\ast$)}, and let $\xi$ be uniform over the odd residues
modulo~$P$. For $1\le a\le a_0$ and $k\in\mathcal K$ put $\kappa=\kappa_a(k)$ and
$\theta_k=N2^{\kappa-k}/(2^{a+4}n_\kappa^2)$, and let $\beta^{(a)}_k$ be the probability that $k$ is not $a$-good in
direction~$+$.
\begin{enumerate}[label=(\alph*)]
\item $\theta_k\le1/2$; $k$ fails to be $a$-good in direction~$-$ also with probability $\beta^{(a)}_k$; and
  $\beta^{(a)}_k\le V^{(a)}_k/(V^{(a)}_k+(1-\theta_k)^2)+2^{\kappa+2-s_0}$.
\item With $\beta_a=\card{\mathcal K}^{-1}\sum_{k\in\mathcal K}\beta^{(a)}_k$, the probability that
  $\Gamma_{a,\pm}(\xi)<\gamma\card{\mathcal K}$ for some $1\le a\le a_0$ and some sign is at most
  $\frac2{1-\gamma}\sum_{a=1}^{a_0}\beta_a$.
\item If Conjecture~S holds for $a=1,\dots,a_0$, there is $m_0=m_0(a_0,\gamma)$ such that for every $m\ge m_0$
  \textup{($\ast$)} holds and, for every $s_0>m$, with probability at least~$1/2$,
  $\Gamma_{a,\pm}(\xi)\ge\gamma\card{\mathcal K}$ for every $1\le a\le a_0$ and both signs.
\end{enumerate}
\end{proposition}

\begin{outline}
(a) As $N\le2n_\kappa$, and $2^{\kappa-k}\le2^{j_0+a}\le2^{a+1}N$ because $2^{j_0}\le2N$, we get
$\theta_k\le1/2$. The threshold of $a$-goodness is $\tau_kM_k$ with $\tau_k=R_k/(2^{a+6}\Pi_\kappa)=\theta_k\mu_k$. So by \cref{lem:dnp-transfer}(c), and by Cantelli's inequality
$\Pr[x\le\mu-\lambda]\le\sigma^2/(\sigma^2+\lambda^2)$ ($\lambda>0$) with $\lambda=(1-\theta_k)\mu_k$, the bound holds; the
sign~$-$ follows from \cref{lem:dnp-transfer}(b).

(b) By~(a), $\card{\mathcal K}-\Gamma_{a,\pm}(\xi)$ has expectation $\beta_a\card{\mathcal K}$, so by Markov's inequality it
exceeds $(1-\gamma)\card{\mathcal K}$ with probability at most $\beta_a/(1-\gamma)$; take a union bound over the $2a_0$
pairs $(a,\pm)$.

(c) For $k\in\mathcal K$, $\kappa\le N/4+j_0+a_0$ and $s_0\ge N$, so $2^{\kappa+2-s_0}\le2^{3+a_0}N2^{-3N/4}$; with
$\theta_k\le1/2$, (a) gives $\beta_a\le4\bar V^{(a)}_m+2^{3+a_0}N2^{-3N/4}$ for every $s_0>m$. Under Conjecture~S this
tends to~$0$ for each $a\le a_0$, and \textup{($\ast$)} holds for all large~$m$. Choose $m_0$ such that for $m\ge m_0$
the size condition holds and $\sum_{a\le a_0}\beta_a\le(1-\gamma)/4$ for every $s_0>m$; then (b) bounds the failure
probability by~$1/2$.
Formally, (a) is \lean{MergeDepth.notGoodProb_le}, with the mean of the profile computed in
\lean{MergeDepth.profileMean_eq} and Cantelli's inequality in \lean{MergeDepth.cantelli_interval}; (b) is
\lean{MergeDepth.prob_some_bad_le}, which holds without the size condition; and (c) is
\lean{MergeDepth.good_shift_prob}, with Conjecture~S as a hypothesis.
\end{outline}

\subsection{Half-blindness and witnesses}\label{sec:dnp-witness}

\begin{definition}[Families and the fine zone]\label{def:dnp-family}
\formal{MergeDepth/DeepNestingBP.lean}{MergeDepth.InFamily,MergeDepth.familyPerPeriod,MergeDepth.HalfBlindPair}
Let $1\le k\le m-2$ and let $h$ be a head. The four \emph{families} of the class $(h,k)$ are the \emph{apex family}, the
apex hops $(\alpha-d,\alpha+d)$ with $1\le d<R_k$ of the level-$k$ V's of head~$h$, and for $j\in\{0,1,2\}$ the
\emph{band-$j$ family}, their valley hops at the cells $v+\sigma x$ with
$R_{k+j+1}<x<R_{k+j}$~\cite[Lemma~5.1(b)]{BoundedNesting2026}. A family is \emph{half-blind} at~$\xi$ if at least half of
its hops with later endpoint in $[0,P)$ are blind. For $0<\eps\le1/8$ the \emph{fine zone} is
$\mathcal F_\eps=\{k\in\Z:m-\lfloor\eps m\rfloor\le k\le m-2\}$.
\end{definition}

The pair is $P$-periodic, so a family has the same fraction of blind hops in every period and on every cover. In Lean, a
family consists of the shifts by multiples of~$P$ of hops that \lean{MergeDepth.family} lists in one period, in the time of
their head, and the fine zone is \lean{MergeDepth.fineZone}.

\begin{lemma}[Families of fine classes]\label{lem:dnp-family}
\formal{MergeDepth/DeepNestingBP.lean}{MergeDepth.familyPerPeriod_subset_class,MergeDepth.class_card_le_family,MergeDepth.hopError_of_inFamily_blind,MergeDepth.familyBlindErr_le_of_thetaServes}
Let $s_0>m$, $1\le k\le m-2$, let $h$ be a head, and let $G$ be a block-pair label modulo~$q$ on the $r$-fold cover.
\begin{enumerate}[label=(\alph*)]
\item Each family of the class $(h,k)$ contains at least a fraction $1/32$ of the hops of the class.
\item All blind hops of the apex family have error $a^{(h)}(k)$, and all blind hops of the band-$j$ family have error
  $b^{(h)}_j(k)$.
\item If $G$ $\theta$-serves level~$k$ for head~$h$ and a family of the class $(h,k)$ is half-blind at~$\xi$, then the
  common error of its blind hops is at most $64\theta q$.
\end{enumerate}
\end{lemma}

\begin{outline}
(a) Per level-$k$~V of head~$h$ the class has $2R_k$ hops~\cite[Lemma~5.1(b)]{BoundedNesting2026}, the apex family
$R_k-1$, and the band-$j$ family ($l=k+j$) $R_{l+1}-1$ if $l<m$ and $R_m-1$ if $l=m$. As $s_0>m$, $R_k\ge2^{m-k+1}\ge8$, and
$R_k\ge16$ if $k+2<m$. So the fractions are at least $\frac12-\frac1{16}$ (apex), $2^{-j-2}-\frac1{2R_k}\ge\frac1{32}$
(band~$j$ with $l<m$; if $j=2$ then $k+2<m$) and $2^{-j-1}-\frac1{2R_k}\ge\frac1{16}$ (band~$j$ with $l=m$).

(b) This is the band decomposition with $F=0$, as $A_V=A_{\lev V}$ and $E_V=E_{\lev V}$ for a block
label~\cite[Definition~5.11]{BoundedNesting2026}.

(c) The class has $4r\Lambda$ hops and total error at most $4r\Lambda\theta q$. By~(a) and periodicity the family has at
least $4r\Lambda/32$ hops, at least half of them blind, each with the common error~$e$ of~(b); so
$e\cdot4r\Lambda/64\le4r\Lambda\theta q$.
Formally, (a)--(c) are first proved for the families listed in one period (\lean{MergeDepth.family_card_ge},
\lean{MergeDepth.pairHopErr_of_blind}, \lean{MergeDepth.familyBlindErr_le_of_served}); in~(b) the potential of the other
head is constant on a blind window, and the tenth paper's stroke recursion gives the increment of the own head. They
transfer to the walk pair as described in \cref{sec:dnp-setting}.
\end{outline}

\begin{lemma}[Half-blindness at a random odd shift]\label{lem:dnp-halfblind}
\formal{MergeDepth/DeepNestingBP.lean}{MergeDepth.sum_nonHalfBlindPair_le,MergeDepth.prob_many_nonHalfBlindPair_le}
Let $0<\eps\le1/8$ and $m<s_0$ with $\eps m\ge4$, and let $\xi$ be uniform over the odd residues modulo~$P$. The expected
number of triples $(h,k,f)$, with $h$ a head, $k\in\mathcal F_\eps$ and $f$ a family of the class $(h,k)$, that are not
half-blind at~$\xi$ is at most $64\eps^2m$. Hence, for $a_0\ge1$ and $\delta=1/(48(a_0+1))$, more than
$\delta\card{\mathcal F_\eps}$ such triples exist with probability at most $6144(a_0+1)\eps$.
\end{lemma}

\begin{outline}
First, a window $(p,u]$ of length~$\ell$ of either head is blind with probability at least
\[
  N^{-1}\sum_{i=1}^m(1-\ell/R_i)^+-2/(NR_m)
\]
over odd shifts; this is~\cite[Proposition~5.17]{BoundedNesting2026} without its level-$0$ term. Indeed, in the time of
the other head, where it is~$U$, each of the $2^{i+2}$ turns of level~$i\ge1$ per period is followed by $R_i$ steps without
a turn, a window that starts in the first $R_i-\ell$ of them is blind, these starts are disjoint, and at least
$((R_i-\ell)^+-1)/2$ of them have the parity that an odd shift requires; as $2^{i+2}R_i=P/N$ and
$\sum_{i=1}^m2^{i+2}/P<2/(NR_m)$, the bound follows. If $\ell\le R_L$ with $L\le m$, the levels $1,\dots,L$ contribute
at least $\sum_{i=1}^L(1-2^{i-L})\ge L-2$ to the sum, and $2/(NR_m)\le1/N$ as $R_m\ge2$.

Let $k\in\mathcal F_\eps$ and $g=m-k$, so $2\le g\le\eps m$. An apex hop of $(h,k)$ has length $2d<R_{k-1}$, so the sum is
at least $k-3$. A valley hop of band~$j$, at level $l=k+j$, has length $(m-l)R_l+2x<(m-l+2)R_l\le R_{l-\lambda}$ with
$\lambda=\lceil\log_2(m-l+2)\rceil\le\lceil\log_2(g+2)\rceil$~\cite[Lemma~5.1(b)]{BoundedNesting2026}, so the sum is at
least $l-\lambda-2$ if $l\ge\lambda$. So every hop of every family of $(h,k)$ fails to be blind with probability at most
$\eps_g=(g+4+\lceil\log_2(g+2)\rceil)/N$ (trivially if $l<\lambda$, as then $\eps_g>1$). Hence the expected non-blind
fraction of a family is at most~$\eps_g$, and by Markov's inequality it fails to be half-blind with probability at
most~$2\eps_g$. With $G=\lfloor\eps m\rfloor$ we have $\eps_g\le4G/N$, as $\lceil\log_2(g+2)\rceil\le g$; summing over the
two heads, the four families and the $G-1$ levels of~$\mathcal F_\eps$ gives at most $64(G-1)G/N\le64\eps^2m$. As
$\card{\mathcal F_\eps}=G-1\ge\eps m/2$, Markov's inequality bounds the probability of more than
$\delta\card{\mathcal F_\eps}$ such triples by $128\eps/\delta=6144(a_0+1)\eps$.
Formally, the blindness bound is \lean{MergeDepth.blindnessBound_holds}, and the lemma is first proved for the families
listed in one period (\lean{MergeDepth.sum_nonHalfBlind_le'}, \lean{MergeDepth.prob_many_nonHalfBlind_le'}); for the walk
pair, half-blindness is the same condition (\lean{MergeDepth.halfBlindPair_iff}).
\end{outline}

\begin{lemma}[Witnesses give unserved levels]\label{lem:dnp-witness}
\formal{MergeDepth/DeepNestingBP.lean}{MergeDepth.witnesses_unserved_pair}
Let $a_0\ge1$, $0\le\gamma\le1$, $0<\eps\le1/8$, $\theta_0=2^{-(a_0+11)}$ and $\delta=1/(48(a_0+1))$, and let $m<s_0$ satisfy
$\eps m\ge4$ and~\textup{($\ast$)}. Let $\xi$ be an odd shift such that
\begin{enumerate}[label=(\roman*)]
\item $\Gamma_{a,\pm}(\xi)\ge\gamma\card{\mathcal K}$ for all $1\le a\le a_0$ and both signs, and
\item at most $\delta\card{\mathcal F_\eps}$ triples $(h,k,f)$ with $k\in\mathcal F_\eps$ are not half-blind at~$\xi$.
\end{enumerate}
Then for all $1\le a_1,a_2\le a_0$, all $K_1,K_2\in\Z$, every $q\ge3$ and every multiple~$r$ of~$q$, at least $c(m+1)-2$
levels are not doubly $\theta_0$-served by the block-pair label of $(B(a_1,K_1),B(a_2,K_2))$ modulo~$q$ on the $r$-fold
cover, where $c=\min\bigl(\gamma/8,\ \eps/(26(a_0+1))\bigr)$.
\end{lemma}

\begin{outline}
Call a triple \emph{bad} if it is not half-blind. As $\eps m\ge4$ and $\eps\le1/8$,
$\card{\mathcal F_\eps}=\lfloor\eps m\rfloor-1\ge\eps m/2$ and $\mathcal F_\eps\subseteq[1,m-2]$. We find levels that are
not $\theta_0$-served for some head, in three cases covering all $q\ge3$.

\emph{Case 1: $q<32$.} Then $2\lfloor\log_2q\rfloor+2\le10$. Split $\mathcal F_\eps$ into
$\lfloor\card{\mathcal F_\eps}/11\rfloor$ disjoint windows of $11$ consecutive levels and a remainder. For a window
$\{k_0,\dots,k_0+10\}$, the run bound with $s=11$, applied to $B(a_1,K_1)$, fails: there are $0\le i\le10$ and a family~$f$
of the class $(1,k_0+i)$, the apex family or the band-$j$ family with $i+j\le10$, whose blind error exceeds $q/8$; these
quantities involve only the levels $k_0,\dots,k_0+10\le m-2$, so they are blind errors of actual hops
(\cref{lem:dnp-family}(b)). Distinct windows give distinct levels, and at most
$\delta\card{\mathcal F_\eps}\le\card{\mathcal F_\eps}/96$ of these triples are bad. If such a triple $(1,k,f)$ is
half-blind and level~$k$ were $\theta_0$-served for head~$1$, \cref{lem:dnp-family}(c) would bound its blind error by
$64\theta_0q<1$, so the error would be~$0$, not more than $q/8$. So at least
$\lfloor\card{\mathcal F_\eps}/11\rfloor-\card{\mathcal F_\eps}/96\ge\card{\mathcal F_\eps}/13-1$ levels are not
$\theta_0$-served for head~$1$.

\emph{Case 2: $q\ge32$ and $4\tilde Q_h+4>q$ for some head~$h$.} The overflow bound for $B(a_h,K_h)$ gives
$\mathrm{rem}<a_h$, and each of the at least $\lfloor\card{\mathcal F_\eps}/(a_h+1)\rfloor\ge\card{\mathcal F_\eps}/(a_0+1)-1$
levels $k\in\mathcal F_\eps$ with $(m-k)\bmod(a_h+1)=\mathrm{rem}$ has a family~$f$ of the class $(h,k)$ (apex, band~$0$ or
band~$1$) with blind error at least $q/16$. If $(h,k,f)$ is half-blind and level~$k$ were $\theta_0$-served for head~$h$,
\cref{lem:dnp-family}(c) would bound this error by $64\theta_0q=2^{-(a_0+5)}q\le q/64$. Discarding the levels of the at
most $\card{\mathcal F_\eps}/(48(a_0+1))$ bad triples leaves at least $\card{\mathcal F_\eps}/(2(a_0+1))-1$ levels that
are not $\theta_0$-served for head~$h$.

\emph{Case 3: $q\ge32$ and $4\tilde Q_h+4\le q$ for both heads.} Suppose that $\tilde Q_1\le\tilde Q_2$. By~(i), at
least $\gamma\card{\mathcal K}$ levels $k\in\mathcal K$ are $a_2$-good in direction~$+$. Fix one and let
$\kappa=\kappa_{a_2}(k)$, a canceller of $B(a_2,\cdot)$ with $\kappa\le m$ and $\Pi_\kappa\le R_k$. Let $\Phi$ be the
potential of~$U$ carrying $B(a_2,K_2)$, so that head~$2$ has potential $\Phi(t+\xi)$ up to a constant, and let
$\alpha\in\mathcal A_\kappa$ be the apex of a level-$\kappa$~V~$V$. By the path sums~\cite[Lemma~5.9]{BoundedNesting2026},
for $0\le o<\Pi_\kappa$ the values $\Phi(\alpha+o)$ and $\Phi(\alpha+\Pi_\kappa+o)$ are a constant~$C_V$ (from the start of
the level-$0$ stroke and the V's containing~$V$), plus $c^{(2)}_\sigb(\kappa)$ or $c^{(2)}_\sigf(\kappa)$ (from~$V$), plus the contribution~$Y(o)$ of the V's
inside the stroke of~$V$ at offset~$o$, which is the same for both strokes because the increments of a block label depend
only on the level~\cite[Lemma~5.12]{BoundedNesting2026}. This is the decomposition in the proof
of~\cite[Lemma~5.20]{BoundedNesting2026}, and it is the hypothesis of Theorem~V$^*$.

Let $\alpha_1\in\mathcal A_k$, $c_0=(\alpha_1+\xi)\bmod P$ and $S=\mathrm{Strad}(c_0,\alpha,\Pi_\kappa,R_k)$. For $d\in S$,
the apex hop $(\alpha_1-d,\alpha_1+d)$ of head~$1$ lies in the class $(1,k)$ with $d<R_k$, its other-head increment is
$\Phi(c_0+d)-\Phi(c_0-d)$ because $\Phi$ has $P$-periodic increments, and so its error is
$e(d)=(A^{(1)}_k+\Phi(c_0+d)-\Phi(c_0-d))\bmod q$ \cite[Lemma~5.10(a)]{BoundedNesting2026}. Theorem~V$^*$ with
$A'=A^{(1)}_k$ and the residue bound give $\sum_{d\in S}e(d)\ge\card S\,q/4$. For fixed~$\alpha_1$ at most one~$\alpha$ has
$S\ne\emptyset$, since $S\ne\emptyset$ puts $c_0$ in the span $[\alpha,\alpha+2\Pi_\kappa)$ and the spans of distinct V's of
one level are disjoint (\lean{DirectedReciprocity.periodVApices_spans_disjoint}~\cite{DirectedReciprocity2026}); and distinct pairs
$(\alpha_1,d)$ give distinct hops. Summing over $\alpha_1$ and~$\alpha$, the hops of the class $(1,k)$ in one period have
total error at least $\frac q4X^+_{k,\kappa}(\xi)\ge\frac q4\cdot2^{k+2}R_k/2^{a_2+7}=q\Lambda/2^{a_2+7}$, as
$2^kR_k=\Lambda$. Errors are $P$-periodic and the class has $4\Lambda$ hops per period, so its mean error is at least
$q/2^{a_2+9}>\theta_0q$: level~$k$ is not $\theta_0$-served for head~$1$. This gives at least
$\gamma\card{\mathcal K}\ge\gamma N/4-1$ levels.

If $\tilde Q_2<\tilde Q_1$, exchange the heads: use the apex hops of head~$2$, centred at the times
$(\alpha_1-\xi)\bmod P$ of head~$1$, the level-$\kappa_{a_1}(k)$ V's of head~$1$, $X^-$, condition~(i) with $a=a_1$, and
the residue bound with the indices $1$ and~$2$ exchanged.

In every case, since $\card{\mathcal F_\eps}\ge\eps m/2$, at least
$\min\bigl(\eps m/(26(a_0+1)),\ \gamma N/4\bigr)-1\ge c(m+1)-2$ levels are not $\theta_0$-served for one of the heads,
hence not doubly $\theta_0$-served.
Formally, the three cases are \lean{MergeDepth.case1_unserved}, \lean{MergeDepth.case2_unserved} and
\lean{MergeDepth.case3_unserved}, through the tenth paper's run bound of Theorem~LU, its overflow band witness and
Theorem~V$^*$, and \lean{MergeDepth.witnesses_unserved} combines them for service measured on the families listed in one
period; a level that is doubly $\theta_0$-served on a cover is served in this sense
(\lean{MergeDepth.ThetaServes.servedOnWitness}).
\end{outline}

\subsection{Block pairs at one odd shift}\label{sec:dnp-BP}

\begin{theorem}[Theorem BP: block pairs at one odd shift, assuming Conjecture~S]\label{thm:BP}
\formal{MergeDepth/DeepNestingBP.lean}{MergeDepth.theoremBP_pair}
Let $a_0\ge1$, and assume that Conjecture~S holds for $a=1,\dots,a_0$. Put $\theta_0=2^{-(a_0+11)}$ and
$c=2^{-16}/(26(a_0+1)^2)$. There is a constant~$C'$, depending only on~$a_0$ (through the rate of convergence in
Conjecture~S), such that for all $m\ge1$ and $s_0>m$ some odd shift~$\xi$ has the following property: for all integers
$1\le a_1,a_2\le a_0$ and $K_1,K_2$, every $q\ge3$ and every multiple~$r$ of~$q$, the block-pair label of
$(B(a_1,K_1),B(a_2,K_2))$ modulo~$q$ on the $r$-fold cover of the pair $x_1=U(s_0,m,1)$, $x_2=x_1(\cdot+\xi)$ doubly
$\theta_0$-serves at most $(1-c)(m+1)+C'$ levels.
\end{theorem}

\begin{outline}
Put $\gamma=1/2$ and $\eps=2^{-16}/(a_0+1)$, so that $c=\min(\gamma/8,\eps/(26(a_0+1)))$. By \cref{prop:S-shift}(c) there
is $m_0$ such that for $m\ge m_0$ and $s_0>m$ the size condition holds and condition~(i) of \cref{lem:dnp-witness} holds at
a uniformly random odd shift with probability at least $1/2$. Let $m_1=\max(m_0,2^{18}(a_0+1))$ and $C'=\max(2,c(m_1+1))$.
For $m\ge m_1$ and $s_0>m$ we have $\eps m\ge4$, and by \cref{lem:dnp-halfblind} condition~(ii) fails with probability at
most $6144(a_0+1)\eps=3/32<1/2$. So both conditions hold at some odd shift~$\xi$, which depends only on $a_0$, $m$
and~$s_0$, and \cref{lem:dnp-witness} leaves at most $(1-c)(m+1)+2\le(1-c)(m+1)+C'$ doubly $\theta_0$-served levels, for
every block pair with $1\le a_1,a_2\le a_0$, every $q\ge3$ and every multiple~$r$ of~$q$. For $m<m_1$ every odd shift has
the property, since $m+1<(1-c)(m+1)+c(m_1+1)\le(1-c)(m+1)+C'$.
Formally, the theorem is first proved for service measured on the families listed in one period
(\lean{MergeDepth.theoremBP'}), and \lean{MergeDepth.theoremBP_pair} follows, since a level that is doubly
$\theta_0$-served on a cover is served in this sense; the hypothesis $q\mid r$ is not used.
\end{outline}

Theorem~BP matters because block labels are how one head alone serves most of its levels (the fraction $a/(a+1)$ in the
single-head regime~\cite[Lemma~5.12]{BoundedNesting2026}). Assuming Conjecture~S, it shows that at one fixed odd shift no
pair of them with block lengths at most~$a_0$, for any modulus, doubly $\theta_0$-serves more than a fraction $1-c$ of the
levels, up to the additive constant~$C'$, as Conjecture~A$''$ predicts for every label of cost at most $n_{\mathrm h}/8$.

\begin{remark}[Scope of Theorem~BP]\label{rem:dnp-scope}
In the notation of Conjecture~A$''$, Theorem~BP gives, for pairs of block labels of length at most~$a_0$ on covers with
$q\mid r$, the constants $\theta_0$, $\eta=1-c$, $C=0$ and~$C'$, without the cost hypothesis; labels with $q\le2$ values
cost at least $n_{\mathrm h}/2$ by Lemma~Zq~\cite[Lemma~2.7]{BoundedNesting2026}, so $q\ge3$ loses nothing. It leaves
open Conjecture~S; block lengths growing with~$q$, such as $a\approx\frac12\log_2q$, where the straddle density $2^{-a}$ in
Case~3 of the proof of \cref{lem:dnp-witness} falls below every fixed~$\theta_0$; and general labels (step~(d)). Even for a
level-uniform label without cancellers the straddle residue $(2A_k+2E_\kappa+2)\bmod q$ can be small, so Case~3 does not
apply. Theorem~BP therefore does not decide $\TCP(2)$; it is evidence for depth-robustness from the labels whose
cancellations are understood.
\end{remark}

\subsection{Partial results on Conjecture S}\label{sec:dnp-partial}

These results are not used in the proof of Theorem~BP. Throughout, $\operatorname{sinc}y=\sin(\pi y)/(\pi y)$ with
$\operatorname{sinc}0=1$, and for $1\le k<\kappa\le m$ and $r\in\Z$
\[
  \varphi_k(r)=\operatorname{sinc}^4\Bigl(\frac{rn_\kappa}{N2^{\kappa+1}}\Bigr)
  \prod_{i=1}^{k-1}\cos^4\Bigl(\frac{\pi rn_i}{N2^i}\Bigr)\prod_{i=k}^{\kappa-1}\cos^2\Bigl(\frac{\pi rn_i}{N2^i}\Bigr).
\]

\begin{proposition}[Exact form of $V_k$]\label{prop:S-fourier}
\formal{MergeDepth/StraddleDensity.lean}{MergeDepth.exact_form_V,MergeDepth.density_W,MergeDepth.sideband_eq_densityW}
Let $1\le k<\kappa\le m$.
\begin{enumerate}[label=(\alph*)]
\item $V_k=\sum_{r\ne0}\varphi_k(r)$.
\item Let $\zeta_1,\dots,\zeta_{\kappa-1},\omega_1,\dots,\omega_4$ be independent, with $\zeta_i+2\sim\mathrm{Bin}(4,\frac12)$
  for $i<k$, $\zeta_i+1\sim\mathrm{Bin}(2,\frac12)$ for $k\le i<\kappa$, and $\omega_j$ uniform on $[-\frac14,\frac14]$, and
  let $W=\sum_{i<\kappa}\zeta_i\hat\Pi_i+Z$ with $Z=\hat\Pi_\kappa(\omega_1+\dots+\omega_4)$. Then $W$ has a continuous
  density~$f_W$, and $1+V_k=N\sum_{l\in\Z}f_W(lN)$.
\item $V_k\ge F_k:=\sum_{t\ne0}\varphi_k(N2^{k-1}t)$, and $F_k=F(n_\kappa,\kappa-k)$, where
  \[
    F(n,D)=\sum_{t\ne0}\operatorname{sinc}^4\bigl(tn2^{-D-2}\bigr)\prod_{e=1}^{D}\cos^2\bigl(\pi t(n+D+1-e)2^{-e}\bigr).
  \]
\end{enumerate}
\end{proposition}

\begin{outline}
(a) With $\Delta_N(y)=\sum_l\Delta(y-lN)$, $x_k(t)=(2^{k-1}\hat\Pi_\kappa)^{-1}\sum_{S,S'}\Delta_N(t+\hat a_S-\hat v_{S'})$,
so the Fourier coefficients $\hat x(r)=N^{-1}\int_0^Nx_k(t)e^{-2\pi irt/N}dt$ are
\[
  \hat x(r)=\frac{\hat\Delta(r/N)}{2^{k-1}\hat\Pi_\kappa N}\sum_Se^{2\pi ir\hat a_S/N}\sum_{S'}e^{-2\pi ir\hat v_{S'}/N},
\]
with $\hat\Delta(\eta)=\int_\R\Delta(y)e^{-2\pi i\eta y}dy=\frac12\hat\Pi_\kappa^2\operatorname{sinc}^2(\hat\Pi_\kappa\eta/2)$
by direct integration. As $\hat a_S$ is a constant plus $\sum_{i\in S}\hat\Pi_i$,
$\abs{\sum_Se^{2\pi ir\hat a_S/N}}=\prod_{i<k}\abs{1+e^{2\pi ir\hat\Pi_i/N}}=2^{k-1}\prod_{i<k}\abs{\cos(\pi rn_i/(N2^i))}$,
and likewise for the valleys with $\kappa-1$ factors. Hence $\abs{\hat x(r)}^2/\hat x(0)^2=\varphi_k(r)$, the factors with
$i<k$ occurring in both sums; $\hat x(0)=\mu_k$, and
Parseval's identity $\operatorname{Var}x_k=\sum_{r\ne0}\abs{\hat x(r)}^2$ gives~(a).

(b) For $\zeta+2\sim\mathrm{Bin}(4,\frac12)$,
$\mathbb Ee^{2\pi iy\zeta}=e^{-4\pi iy}((1+e^{2\pi iy})/2)^4=\cos^4(\pi y)$; for $\zeta+1\sim\mathrm{Bin}(2,\frac12)$ it is
$\cos^2(\pi y)$; for $\omega$ uniform on $[-\frac14,\frac14]$ it is $\operatorname{sinc}(y/2)$. By independence,
$\psi(y)=\mathbb Ee^{2\pi iyW}$ is the product of these factors at $y\hat\Pi_i$ and $y\hat\Pi_\kappa$, and
$\psi(r/N)=\varphi_k(r)$ for $r\in\Z$. Let $f_W$ be the mixture, with the probabilities of the values of
$(\zeta_1,\dots,\zeta_{\kappa-1})$ as weights, of the translates by $\sum_i\zeta_i\hat\Pi_i$ of the density of~$Z$ (a
constant times the convolution square of~$\Delta$). It is continuous, nonnegative and compactly supported, and its Fourier
transform is~$\psi$, so it is the density of~$W$. Poisson summation for $y\mapsto f_W(Ny)$, whose Fourier transform at
$r\in\Z$ is $\psi(r/N)/N=\varphi_k(r)/N$ ($\varphi_k$ is even) and is summable as
$\operatorname{sinc}^4y\le\min(1,(\pi y)^{-4})$, gives $N\sum_lf_W(lN)=\sum_r\varphi_k(r)=1+V_k$ by~(a).

(c) All terms in~(a) are nonnegative. At $r=N2^{k-1}t$ the factors with $i<k$ equal~$1$, the factor at $k\le i<\kappa$ is
$\cos^2(\pi tn_i2^{k-1-i})$, and the sinc factor is $\operatorname{sinc}^4(tn_\kappa2^{k-\kappa-2})$; with $D=\kappa-k$,
$e=i-k+1$ and $n_i=n_\kappa+D+1-e$, this is the term~$t$ of $F(n_\kappa,D)$.
Formally, $\varphi_k$ is \lean{MergeDepth.straddlePhi}; (a) is \lean{MergeDepth.hasSum_straddlePhi_ne_zero}, through
Mathlib's Parseval identity on the circle of length~$N$, and holds for $1\le k\le\kappa$ (for every $\kappa\ge1$ in the
general form \lean{MergeDepth.hasSum_phiGen}); (c) is \lean{MergeDepth.fineFloorK_le_relVariance} and
\lean{MergeDepth.fineFloorK_eq}, and only the identity $F_k=F(n_\kappa,\kappa-k)$ uses $\kappa\le m$;
\lean{MergeDepth.exact_form_V} collects (a) and~(c). Part~(b) is \lean{MergeDepth.density_W}, for independent random
variables with these laws on any probability space (they exist, \lean{MergeDepth.density_W_model_exists}); it needs only
$1\le\kappa$ and $k\le\kappa$, the law of~$W$ is identified through characteristic functions, and the identity
$\mathbb E\exp(2\pi irW/N)=\varphi_k(r)$ is \lean{MergeDepth.charFun_lawW_eq_straddlePhi}.
\end{outline}

So $1+V_k$ is $N$ times the density at~$0$ of a random sum of the digits $\hat\Pi_i$ modulo~$N$, smoothed at the scale
$\hat\Pi_\kappa$. For $k\in\mathcal K$ and $\kappa=\kappa_a(k)$, splitting the frequencies $r\ne0$ into the multiples of
$N2^{k-1}$, the \emph{semi-local} ones $r=N2^{k-h}s$ ($s$ odd, $2\le h\le k$) and the \emph{sidebands} $r\notin N\Z$ gives
$V_k=F_k+\sum_{h=2}^kSL^{[h]}_k+SB_k$ with $SL^{[h]}_k=\sum_{s\text{ odd}}\varphi_k(N2^{k-h}s)$ and
$SB_k=\sum_{r\notin N\Z}\varphi_k(r)$, all nonnegative (formally \lean{MergeDepth.relVariance_decomp}, for $1\le k\le\kappa$).
The argument of~(b) with period~$1$ gives $\sum_w\varphi_k(Nw)=\sum_lf_W(l)$, so $SB_k=N\sum_lf_W(lN)-\sum_lf_W(l)$
(\lean{MergeDepth.sideband_eq_densityW}) measures how far the law of~$W$ modulo~$N$ is from invariance under integer
translations.

\begin{proposition}[Classes with a large fine floor]\label{prop:S-resonance}
\formal{MergeDepth/StraddleResonance.lean}{MergeDepth.resonance_class}
Let $a=1$, let $q\ge6$ be even, let $k\ge1$ with $3k\le2^q+4$, and let $m=2^q+k+3$. Then $k\in\mathcal K$,
$\kappa_1(k)=k+q+2$, $n_{\kappa_1(k)}=2^q-q+2$, and $V^{(1)}_k\ge F_k>0.585$.
\end{proposition}

\begin{outline}
Here $N=2^q+k+4$ and $k+4\le(2^q+16)/3<2^q$, so $j_0=q+1$, and $4k\le N$, so $k\in\mathcal K$. As $m-k-j_0=2^q-q+2$ is even,
the least $\kappa\ge k+q+1$ with $m-\kappa$ odd is $k+q+2$, so $n_\kappa=2^q-q+2$ and $D=q+2$. By
\cref{prop:S-fourier}(c), $V^{(1)}_k\ge F_k=F(n_\kappa,q+2)$, whose cosine factors are $\cos^2(\pi t(2^q+5-e)2^{-e})$ for
$1\le e\le q+2$. Keep only the two equal terms $t=\pm4$. The sinc factor is $\operatorname{sinc}^4(y)$ with
$y=(2^q-q+2)2^{-q-2}\in(0,\frac14]$, so it is at least $\operatorname{sinc}^4(\frac14)=(2\sqrt2/\pi)^4>0.65702$. For $t=4$
the factor of index~$e$ is $\cos^2(\pi(2^q+5-e)2^{2-e})$: it is~$1$ for $e\in\{1,2,3,5\}$; $\frac12$ for $e=4$; and
$\cos^2(\pi(e-5)2^{2-e})$ for $6\le e\le q+2$, since $2^{q+2-e}$ is an integer. These factors lie in $[0,1]$, so their
product is at least
$\prod_{e\ge6}\cos^2(\pi(e-5)2^{2-e})\ge\cos^4(\pi/16)\cos^2(3\pi/64)(1-\pi^2\sum_{j\ge4}j^24^{-j-3})>0.89142$, using the
factors $e=6,7,8$ exactly, $\cos^2y\ge1-y^2$ and $\prod(1-y_e)\ge1-\sum y_e$ for $e\ge9$, and
$\sum_{j\ge4}j^24^{-j}=\frac{20}{27}-\frac{41}{64}$. Hence
$F_k\ge2\cdot0.65702\cdot\frac12\cdot0.89142>0.585$.
Formally, the numerical bounds come from the exact value of $\cos(\pi/8)$, the bound $\cos y\ge1-y^2/2$ and Mathlib's
bounds on~$\pi$. The bound on~$F_k$ does not use that $q$ is even; evenness gives $\kappa_1(k)=k+q+2$.
\end{outline}

So $V^{(1)}_k\not\to0$ along classes $k\to\infty$ (this family has at most one class per~$m$), and variance bounds for
single classes cannot show that every coarse class is good; Conjecture~S averages over classes, as \cref{prop:S-shift}(b)
allows.

\paragraph{Arithmetic of $\kappa_a$.} Let $a\ge1$ with $4(j_0+a)\le N$, and for $k\in\mathcal K$ let
$\rho(k)=(m-k-j_0)\bmod(a+1)$ and $\ell_k=\lfloor(m-k-j_0+1)/(a+1)\rfloor$. Then $\kappa_a(k)=k+D(k)$ with
$D(k)=j_0+((\rho(k)+1)\bmod(a+1))\in[j_0,j_0+a]$, which depends only on~$\rho(k)$, and $n_{\kappa_a(k)}=(a+1)\ell_k$. The
map $k\mapsto(\rho(k),\ell_k)$ is injective on~$\mathcal K$, and the $\ell_k$ lie in an interval~$I$ of integers with
$(a+1)\card I\le\lfloor N/4\rfloor+2a\le N$. Moreover $2^{D(k)}\ge2^{j_0}>N$, and $N\le2n_{\kappa_a(k)}$ and
$2^{j_0}\le2N$ give $2^{D(k)}\le2^{a+2}n_{\kappa_a(k)}$, that is, $n_{\kappa_a(k)}2^{-D(k)-2}\ge2^{-a-4}$.

\begin{lemma}[Dyadic products along a progression]\label{lem:S-progression}
\formal{MergeDepth/StraddleProgression.lean}{MergeDepth.window_sum_prod_cos_sq,MergeDepth.progression_avg_le_strong}
Let $j\ge0$.
\begin{enumerate}[label=(\alph*)]
\item If $u$ is an odd integer and $\gamma_1,\dots,\gamma_j$ are real, then for every integer~$\ell_0$
  \[
    \sum_{\ell=\ell_0}^{\ell_0+2^j-1}\ \prod_{s=1}^j\cos^2\bigl(\pi(u\ell2^{-s}+\gamma_s)\bigr)=1.
  \]
\item Let $a+1=2^wo$ and $t=2^vt'$ with $o$, $t'$ odd and $w,v\ge0$, let $c_e$ be real numbers for $v+w<e\le v+w+j$, and
  let $I$ be a set of $L\ge1$ consecutive integers. Then
  \[
    \frac1L\sum_{\ell\in I}\ \prod_{e=v+w+1}^{v+w+j}\cos^2\bigl(\pi t((a+1)\ell+c_e)2^{-e}\bigr)\ \le\ 2^{-j}+\frac{1-2^{-j}}L.
  \]
\end{enumerate}
\end{lemma}

\begin{outline}
(a) For $j=0$ the sum has the single term~$1$. For $j\ge1$, replacing $\ell$ by $\ell+2^{j-1}$ shifts the arguments with
$s<j$ by the integer multiples $u2^{j-1-s}\pi$ of~$\pi$ and the argument with $s=j$ by $\pi u/2$, which turns $\cos^2$ into
$\sin^2$ as $u$ is odd. So pairing $\ell$ with $\ell+2^{j-1}$ for $\ell_0\le\ell<\ell_0+2^{j-1}$ turns the last factor
into $\cos^2+\sin^2=1$ and reduces the sum to the same sum for $j-1$.

(b) Put $u=t'o$, which is odd, and $e=v+w+s$; then $t(a+1)2^{-e}=u2^{-s}$, and the product is
$\prod_{s=1}^j\cos^2(\pi(u\ell2^{-s}+\gamma_s))$ with $\gamma_s=tc_e2^{-e}$. Split~$I$ into $\lfloor L/2^j\rfloor$ runs of
$2^j$ consecutive integers and a rest of fewer than $2^j$. By~(a) each run contributes exactly~$1$, and the rest at
most~$1$, as the terms are nonnegative; so the sum is at most $\lceil L/2^j\rceil\le L2^{-j}+1-2^{-j}$.
Formally, (a) is \lean{MergeDepth.window_sum_prod_cos_sq} and (b) is \lean{MergeDepth.progression_avg_le_strong}.
\end{outline}

\begin{proposition}[Fine floors and near semi-local frequencies]\label{prop:S-local}
\formal{MergeDepth/StraddleLocal.lean}{MergeDepth.sum_fineFloorK_le,MergeDepth.sum_semiLocal_le,MergeDepth.relVariance_decomp,MergeDepth.conjectureS_iff_meanNonLocalGen}
Let $a\ge1$, for $k\in\mathcal K$ let $\kappa=\kappa_a(k)$, and put $C_a=(a+1)\bigl(3(2^{a+4}/\pi)^2+8\cdot2^{a+4}/(3\pi)\bigr)$.
\begin{enumerate}[label=(\alph*)]
\item If $4(j_0+a)\le N$, then $\sum_{k\in\mathcal K}F_k\le C_a$.
\item If $4(j_0+a)\le N$, then $\sum_{k\in\mathcal K}\sum_{h=2}^{\min(H,k)}SL^{[h]}_k\le(a+1)2^{a+H+8}/(3\pi)$ for every
  integer $H\ge0$.
\item Let $H=H(m)\ge0$ be integers with $2^H/N\to0$ as $m\to\infty$, and put
  $NL_k=V_k-F_k-\sum_{h=2}^{\min(H,k)}SL^{[h]}_k$, which is $SB_k+\sum_{h=\max(2,H+1)}^kSL^{[h]}_k$. Then Conjecture~S holds
  for~$a$ if and only if $\card{\mathcal K}^{-1}\sum_{k\in\mathcal K}NL_k\to0$ as $m\to\infty$.
\end{enumerate}
\end{proposition}

\begin{outline}
Write $a+1=2^wo$ with $o$ odd, and let $I$ be the interval of the arithmetic above, so that $2^w\card I\le(a+1)\card I\le N$.

(a) By the arithmetic above, $F_k=F((a+1)\ell_k,D(k))$ with $n_\kappa2^{-D(k)-2}\ge2^{-a-4}$. As
$\abs{\operatorname{sinc}y}\le\min(1,1/(\pi\abs y))$, the sinc factor of the term~$t$ of~$F_k$ is at most
$\mathrm{sb}(t)=\min(1,(\pi\abs t2^{-a-4})^{-4})$, so $F_k\le M(\ell_k,D(k))$ with
$M(\ell,D)=\sum_{t\ne0}\mathrm{sb}(t)\prod_{e=1}^D\cos^2(\pi t((a+1)\ell+D+1-e)2^{-e})\ge0$. As $k\mapsto(\rho(k),\ell_k)$ is
injective, $\sum_{k\in\mathcal K}F_k\le\sum_{\rho\le a}\sum_{\ell\in I}M(\ell,D(\rho))$. Fix $D=D(\rho)$ and $t\ne0$, and
write $t=2^vt'$ with $t'$ odd. If $v+w<D$, drop the cosine factors with $e\le v+w$, which lie in $[0,1]$, and apply
\cref{lem:S-progression}(b) with $j=D-v-w$ and $c_e=D+1-e$: the sum over $\ell\in I$ of the cosine products is at most
$\card I2^{v+w-D}+1$. If $v+w\ge D$, it is at most $\card I\le\card I2^{v+w-D}$. As $2^v\le\abs t$ and $2^w\card I\le N<2^D$,
it is at most $\abs t+1$ in both cases. So $\sum_{\ell\in I}M(\ell,D)\le\sum_t(\abs t+1)\,\mathrm{sb}(t)$, and telescoping
bounds give $\sum_t\abs t\,\mathrm{sb}(t)\le3(2^{a+4}/\pi)^2$ and $\sum_t\mathrm{sb}(t)\le8\cdot2^{a+4}/(3\pi)$.

(b) Fix $2\le h\le H$, and for $k\in\mathcal K$ with $k\ge h$ put $E=D(k)+h-1$. At $r=N2^{k-h}s$ with $s$ odd, the factors
of~$\varphi_k$ with $i\le k-h$ equal~$1$. Bound the $\cos^4$ factors by~$\cos^2$, keep only the factors with $i=k-h+l$,
$w<l\le E$, which are $\cos^2(\pi s(n_\kappa+E+1-l)2^{-l})$, and bound the sinc factor
$\operatorname{sinc}^4(sn_\kappa2^{-D(k)-h-1})$ by $\min(1,(\pi\abs s2^{-a-3-h})^{-4})$. As in~(a), pass to the pairs
$(\rho,\ell)$: \cref{lem:S-progression}(b) with $t=s$, $v=0$ and $j=E-w$ bounds the sum over $\ell\in I$ of the kept products
by $\card I2^{w-E}+1\le2$, as $2^E\ge2^{j_0}>N$, and the sinc bounds sum to at most $8\cdot2^{a+3+h}/(3\pi)$. So
$\sum_{k\in\mathcal K,\,k\ge h}SL^{[h]}_k\le(a+1)2^{a+7+h}/(3\pi)$, and summing over $2\le h\le H$ gives the bound.

(c) The second form of $NL_k$ follows from the decomposition $V_k=SB_k+F_k+\sum_{h=2}^kSL^{[h]}_k$. For all large~$m$
the size condition holds, and then (a) and~(b) give
\[
  0\le\bar V^{(a)}_m-\frac1{\card{\mathcal K}}\sum_{k\in\mathcal K}NL_k
  =\frac1{\card{\mathcal K}}\sum_{k\in\mathcal K}\Bigl(F_k+\sum_{h=2}^{\min(H,k)}SL^{[h]}_k\Bigr)
  \le\frac{C_a+(a+1)2^{a+H+8}/(3\pi)}{\card{\mathcal K}},
\]
which tends to~$0$ as $\card{\mathcal K}=\lfloor N/4\rfloor$ and $2^H/N\to0$; so the two averages tend to~$0$ together.
Formally, (a), (b) and~(c) are \lean{MergeDepth.sum_fineFloorK_le}, \lean{MergeDepth.sum_semiLocal_le} and
\lean{MergeDepth.conjectureS_iff_meanNonLocalGen}, and the decomposition is \lean{MergeDepth.relVariance_decomp}; for
$H=\lfloor\frac14\log_2N\rfloor$, (c) is also \lean{MergeDepth.conjectureS_iff_meanNonLocal}.
\end{outline}

So the fine floors and the semi-local frequencies with $2^h=o(N)$, for instance
$h\le\lfloor\log_2N\rfloor-\lceil\log_2\log_2N\rceil$, contribute $o(1)$ to~$\bar V^{(a)}_m$, and Conjecture~S is a statement
about the sidebands and the semi-local frequencies with $2^h$ at least of order~$N$. By~(a), at most $2C_a$ coarse classes have
$F_k\ge1/2$, for every $m$ with $4(j_0+a)\le N$; the classes of \cref{prop:S-resonance} are of this kind.

\begin{remark}[The carry chain]\label{rem:S-carry}
Proof sketch only: a representation for the non-local part. Let $b_1,\dots,b_{\kappa-1}$ be the binary digits of
$\sum_{i<\kappa}\zeta_i2^{-i}\bmod1$, and define carries by $\varrho_\kappa=0$ and $\varrho_i=(\zeta_i+\varrho_{i+1}-b_i)/2$
for $i=\kappa-1,\dots,1$. Each $\zeta_i$ is odd with probability~$\frac12$, so the $b_i$ are independent fair bits, and
given them the carries, which stay in $\{-2,\dots,2\}$, form a Markov chain run downwards from level $\kappa-1$.
Telescoping gives $\sum_i\zeta_i\hat\Pi_i\equiv\sum_ib_i\hat\Pi_i-2\varrho\pmod N$ with $\varrho=\sum_i\varrho_i2^{-i}$,
so $1+V_k=N\,\mathbb E\bigl[\sum_lf_Z\bigl(\sum_ib_i\hat\Pi_i-2\varrho-lN\bigr)\bigr]$, where $f_Z$ is the density of~$Z$
and $\sum_ib_i\hat\Pi_i$ is a uniformly random valley position of the comb. The difficulty is the correlation of
the carries with the leading bits, the mixed-radix carry problem that also blocked the fixed-shift approach
of~\cite[Section~5.6]{BoundedNesting2026}: the carries forget distant bits quickly (\cref{rem:S-numerics}), but a proof
needs near-invariance of the law of~$W$ modulo~$N$ under integer translations, which $SB_k$ measures.
\end{remark}

\begin{remark}[Numerical: Conjecture S]\label{rem:S-numerics}
The values below were computed from \cref{prop:S-fourier}(b) by propagating the exact law of the partial sums
$\sum_{i\le j}\zeta_i\hat\Pi_i$ level by level ($O(\kappa N)$ floating-point operations per class, no sampling); where
both were run, this agrees with direct summation of \cref{prop:S-fourier}(a) to four or five decimals.
\begin{enumerate}[label=(\arabic*)]
\item For $a=1$, $\bar V^{(1)}_m=0.268$, $0.134$, $0.061$, $0.0345$, $0.0189$, $0.0114$ at $m=127$, $255$, $511$, $1023$,
  $2047$, $4095$; at the eleven values of~$m$ computed between $127$ and $4095$, $m\bar V^{(1)}_m$ lies between $30$
  and~$47$, consistent with $\bar V^{(1)}_m=O(\log m/m)$. For $a=2$ and $a=3$, $\bar V^{(a)}_m=0.126$, $0.058$, $0.031$,
  $0.021$ and $0.139$, $0.064$, $0.030$, $0.018$ at $m=255$, $511$, $1023$, $2047$.
\item The sidebands of the first classes carry most of $\bar V$: for $a=1$ and $m\ge511$ the averages of $F_k$ and
  $\sum_hSL^{[h]}_k$ are at most $0.0017$ and $0.0044$, while $\sum_kSB_k\in[6.5,10.5]$ for all eleven~$m$; at $m=2047$,
  $SB_k\approx1.68$, $0.28$, $0.024$, $0.0017$, $10^{-5}$ at $k=1$, $8$, $16$, $24$, $40$.
\item By concavity of $v\mapsto v/(v+s)$, \cref{prop:S-shift}(a) gives
  $\beta_a\le\bar V^{(a)}_m/(\bar V^{(a)}_m+(1-\theta^*)^2)+2^{3+a_0}N2^{-3N/4}$ with $\theta^*=\max_k\theta_k\approx0.23$, so
  $\beta_1\le0.055$, $0.031$, $0.019$ at $m=1023$, $2047$, $4095$; at $m=1023$, $2047$ the event of \cref{prop:S-shift}(c)
  has probability at least $\frac12$ for all $s_0>m$, for $(a_0,\gamma)=(2,\frac12)$ ($\beta_1+\beta_2\le0.105$, $0.065$)
  and $(3,\frac14)$ ($\beta_1+\beta_2+\beta_3\le0.154$, $0.094$).
\item At the classes of \cref{prop:S-resonance}, $F_k=0.641$, $0.604$, $0.593$, $0.590$ for $q=6$, $8$, $10$, $12$; at
  $(m,k)=(85,18)$ and $(89,22)$ the class is not $1$-good at about $24$--$29\%$ of sampled shifts.
\item By exact enumeration, two copies of the carry chain of \cref{rem:S-carry} driven by the same bits coalesce within
  two levels $i<k$ with probability at least $5/16$, whatever the bits and carries; so if $p+2j<k$, the bits $b_i$ with
  $i>p+2j$ move the law of the carries at the levels $i\le p$ by at most $(11/16)^j$ in total variation.
\end{enumerate}
\end{remark}
